
\section{Full experimental protocol}
\label{app:full-experiments}

\subsection{Experimental questions}
\label{subsec:experimental-questions}

The experiments are designed to answer four predeclared questions:

\begin{enumerate}
\item Does StableMax-PPO attain a competitive held-out robust
\(\mathbb E[\max_{j\le32}R_j]\) objective relative to
mean optimization and a direct scalar max@32 policy gradient, while
reducing held-out KL?
\item Are gains preserved under the nominal exact ordinal maximum and the
top-rating probability \(\Pr\{\exists j:R_j=4\}\)?
\item Do nominal expected rating, independent Quality RM score, and
reference-policy KL remain within their predeclared safeguards?
\item Do independent generated-output proxies and existing human-labeled
preference diagnostics provide evidence against target-RM exploitation?
\end{enumerate}

No result is used to remove a comparator, redefine a threshold, or select a more
favorable protocol.

\subsection{Protocol audit and scientific interpretation}
\label{subsec:gate-outcomes}

The artifact verifier passed for seeds 314, 2718, and 1618.  The terminal
record preserves every predeclared proxy, non-inferiority, safeguard, and KL
outcome for auditability. We interpret these fields as protocol diagnostics
that characterize the measured operating points, not as a universal method
ranking or a proof that reward-model exploitation is absent.

The diagnostic vector is read as a decision record rather than a single
leaderboard score.  A method can satisfy KL and quality safeguards while
losing a max comparison, or improve the target reward while degrading an
independent proxy.  Reporting the vector preserves these distinctions and
prevents a favorable endpoint from hiding a failure mode that the method was
designed to measure.

The aggregate terminal record reports \texttt{scientific\_success=false}:
the mean and independent-quality floors and the held-out KL comparisons against
Mean-PPO and GRPO pass, but the declared robust-maximum gates and Qwen
proxy gate do not all pass. In particular, the paired robust-maximum difference
from Mean-PPO is 0.0339 with a 95\% interval of $[-0.0686,0.1304]$.
The corresponding intervals against GRPO and Scalar max-PPO are
$[-0.1336,-0.0609]$ and $[-0.1024,-0.0101]$. Passing artifact verification
therefore establishes protocol consistency, not scientific success.

The recorded runs fix \(\varepsilon_{\rm run}=0.04680826120821986\)
using the earlier calibration rule with \(\alpha_{\rm cal}=0.05\).
As explained in Appendix~\ref{app:theory}, this is a fixed ambiguity-set
parameter, not a certified population calibration radius. We retain all
recorded outcomes and do not substitute the corrected formula into completed
runs. The released summaries record the radius; the original cluster-level
calibration data and report remain external, so their sample count and split
independence cannot be reconstructed from this release.

\subsection{Models, data, and split discipline}
\label{subsec:models-data}

The policy is Qwen2.5-1.5B-Instruct
\citep{qwen2024qwen25}.  The frozen Moment RM is a Qwen2.5-14B ordinal
model trained on HelpSteer2 repeated helpfulness ratings, with the
decontaminated preference replay recorded in its training manifest
\citep{wang2024helpsteer2}.  It exports
\(p_{\omega,0:4}\) directly; exposing the category probabilities
does not retrain or alter its accepted checkpoint.  An independently trained
Qwen2.5-14B scalar Quality RM, fit separately on disjoint preference data,
supplies \(Q_\eta\) and has no
variance head.  Its frozen 4,096-pair confirmation accuracy is
\(94.60\%\) with a Wilson lower 95\% bound of
\(93.87\%\).

All policy, reward-model, evaluator, and lockbox splits are content-hashed.
Before partitioning RewardBench, we remove prompt and response overlap with
the policy and Quality-RM data \citep{lambert2024rewardbench}.  The remaining
calibration, policy-diagnostic, and reserve partitions are also mutually
prompt- and response-disjoint.  They contain 1,024, 1,024, and 427 pairs,
respectively.  The sealed Skywork diagnostic contains 4,096 pairs and is
excluded from Quality-RM training and selection
\citep{liu2024skyworkreward}.

The Moment RM and policy prompts remain drawn from the same broad HelpSteer2
domain.  Disjoint rows prevent direct leakage but do not establish
out-of-domain generalization; this is treated as a limitation rather than
hidden by the split terminology.

\subsection{Frozen training protocol}
\label{subsec:experimental-protocol}

The primary response budget is \(N=32\).  Each trainable confirmatory method uses
three seeds \(\{314,2718,1618\}\), 300 policy steps, two prompts
per step, 32 responses per prompt, learning rate \(2\times10^{-5}\),
one PPO epoch, clip \(0.2\), LoRA rank 8 with scale 32, sampling
temperature 1, \(\mathrm{top}_p=1\), maximum prompt length 384,
and maximum response length 64. Seed \(42\) is used only for development and real-GPU smoke tests and is
excluded from all confirmatory intervals, gates, and aggregate claims. The target KL is \(0.02\) and the
hard rollback level is \(0.04\).  All methods share the same prompt
stream, generation budget, optimizer class, policy architecture, and
checkpoint-selection rule.

The non-inferiority floors are frozen before any trainable policy outcome is
observed.  The reference sample contains 2,048 base-policy responses from 512
independent policy-training prompts, with four responses per prompt.  For
\(Z\in\{\mu_\omega,Q_\eta\}\), let
\(\bar Z_g\) be the four-response prompt mean and let
\(L_Z\) be the fifth percentile of 10,000 bootstrap means obtained by
resampling prompts.  The frozen floor is \(z_0=L_Z-0.1s_Z\),
where \(s_Z\) is the response-level standard deviation in the same
reference sample.  The bootstrap unit, one-sided level, 10,000 draws, sample
counts, generation seed, bootstrap seed, and practical margin are all
predeclared.  A separate prompt-disjoint sample of 512 responses from 128
prompts reports standardized reference-policy shift.  It is diagnostic only
and cannot accept or reject training.

\subsection{Comparators and ablations}
\label{subsec:comparators}

The ten trainable methods are:

\begin{enumerate}
\item \textbf{StableMax-PPO}: the complete method in
Algorithm~\ref{alg:stablemax}.
\item \textbf{Mean-PPO + Q}: PPO on
\(\mu_\omega\) with the same mean, Quality, and KL protocol.
\item \textbf{GRPO + Q}: group-relative standardized
\(\mu_\omega\) credit with matched safeguards
\citep{shao2024deepseekmath,guo2025deepseekr1}.
\item \textbf{Scalar Max@32-PPO + Q}: exact leave-one-out scalar
max credit
\( [\mu_i-\max_{j\ne i}\mu_j]_+ \)
with matched safeguards
\citep{tang2025inference,walder2025pkpo,bagirov2025best}.
\item \textbf{Entropic-PPO + Q}: PPO on
\(\mathrm{CE}_{\tau}=-\tau\log\sum_r p_{\omega,r}e^{-r/\tau}\),
using \(\tau=1\).
\item \textbf{Nominal Ordinal EV-PPO + Q}: exact
\(J_{32}^{\mathrm{nom}}\) without the CDF ambiguity set.
\item \textbf{StableMax-PPO without mean floor}: ablates nominal-mean
non-inferiority.
\item \textbf{StableMax-PPO without Quality RM}: ablates the independent
quality safeguard and is excluded from constrained-feasibility claims.
\item \textbf{Gaussian EV-PPO + Q}: the previous heteroscedastic
Gaussian expected-maximum objective, retained as a model-assumption ablation.
\item \textbf{Top-4 PPO + Q}: exact optimization of
\(\Pr\{\exists j:R_j=4\}\), isolating the highest
threshold from the complete ordinal maximum.
\end{enumerate}

An eleventh, evaluation-only \textbf{Best-of-32} row samples 32 responses
from the untrained base policy and ranks them by
\(\mu_\omega\).  It is an inference-time reference, not a
compute-matched training baseline; its generation and scoring cost is reported
separately.

The comparator set is organized to separate three changes that are otherwise
easy to conflate: the deployment functional (mean, scalar maximum, top-rating,
or full ordinal maximum), the CDF robustification, and the stabilizers.
The nominal ordinal and Gaussian variants hold the policy optimizer fixed while
changing the evaluator assumption; the no-floor variants hold the target fixed
while removing one safeguard.  This is why the ablations support a controlled interpretation of the observed
operating points without claiming that any single component is universally
beneficial.

\begin{table}[t]
\caption{Complete aggregate table. Values are three-seed robust expected
maximum, with 95\% confidence intervals, held-out KL, and rollback counts.}
\label{tab:full-aggregate}
\centering\footnotesize
\begin{tabular}{lcccc}
\toprule
Method & Robust max & 95\% CI & Held-out KL & Rollbacks/900\\
\midrule
StableMax-PPO & 3.458 & [3.411, 3.506] & 0.034 & 279 (31.0\%)\\
Mean-PPO & 3.424 & [3.354, 3.517] & 0.070 & 527 (58.6\%)\\
GRPO & 3.555 & [3.530, 3.578] & 0.094 & 741 (82.3\%)\\
Scalar max-PPO & 3.517 & [3.476, 3.558] & 0.054 & 285 (31.7\%)\\
Entropic PPO & 3.400 & [3.353, 3.447] & 0.037 & 302 (33.6\%)\\
Nominal EV-PPO & 3.421 & [3.377, 3.467] & 0.028 & 121 (13.4\%)\\
StableMax-PPO (no mean floor) & 3.508 & [3.426, 3.574] & 0.048 & 230 (25.6\%)\\
StableMax-PPO (no quality floor) & 3.467 & [3.402, 3.530] & 0.048 & 329 (36.6\%)\\
Gaussian EV-PPO & 3.558 & [3.533, 3.581] & 0.061 & 367 (40.8\%)\\
Top-4 PPO & 3.418 & [3.353, 3.499] & 0.042 & 277 (30.8\%)\\
Best-of-32 (base) & 3.086 & [3.045, 3.125] & -- & --\\
\bottomrule
\end{tabular}
\end{table}

\begin{table}[t]
\centering
\footnotesize
\caption{Per-seed primary endpoints at $N=32$. Each cell reports robust
expected maximum / held-out KL on the same 512-prompt evaluation split. The KL
values are the common-base-trajectory means from the held-out evaluator, matching
the aggregate endpoint in Table~\ref{tab:full-aggregate}.}
\label{tab:per-seed-primary}
\begin{tabular}{lccc}
\toprule
Method & Seed 314 & Seed 2718 & Seed 1618 \\
\midrule
StableMax-PPO & 3.503 / 0.057 & 3.452 / 0.024 & 3.420 / 0.022 \\
Mean-PPO & 3.369 / 0.070 & 3.523 / 0.074 & 3.381 / 0.066 \\
GRPO & 3.563 / 0.094 & 3.547 / 0.097 & 3.553 / 0.091 \\
Scalar max-PPO & 3.511 / 0.052 & 3.555 / 0.080 & 3.484 / 0.031 \\
\bottomrule
\end{tabular}
\end{table}

The per-seed view makes the operating-point claim testable without relying on
the aggregate interval alone: StableMax-PPO has lower held-out KL than Mean-PPO
and GRPO in every seed, and lower KL than Scalar max-PPO in two of the three
seeds (2718 and 1618).  Its robust maximum remains competitive across the same paired prompts.
The robust values are the checked-in per-seed CSV records and the KL values are
the held-out common-base means; all entries are rounded only for display, with
no additional aggregation or seed selection performed here.

\subsection{Ablation and stability analysis}
\label{subsec:ablation-stability}

The component rows in Table~\ref{tab:full-aggregate} isolate the operating
point controls.  Removing the mean floor increases the aggregate robust maximum
from \(3.458\) to \(3.508\), but also raises held-out KL from \(0.034\) to
\(0.048\).  Removing the independent Quality floor gives
\((3.467,0.048)\), while removing CDF robustification gives
\((3.421,0.028)\), in (robust maximum, KL) coordinates.  The three-seed
intervals overlap and the per-seed effects are heterogeneous, so these are
sensitivity comparisons rather than isolated causal effects. CDF-envelope
removal changes the evaluator-defined target; floor removal changes the
accepted-update trajectory.  The records therefore support a controlled operating-point interpretation,
not a claim that any safeguard is a universally improving reward term.

\subsection{Evaluation endpoints}
\label{subsec:evaluation-endpoints}

Every method is evaluated on 512 frozen test prompts with newly generated
groups of 32 responses.  The primary paired prompt-level endpoint is
\(h_{32}^{\mathrm{rob}}\), exactly matching the StableMax-PPO
training objective.  Secondary model-based endpoints are
\(h_{32}^{\mathrm{nom}}\),
\(U_{32}^{\mathrm{nom}}\), candidate-level
\(\mu_\omega\), ordinal variance,
\(p_{\omega,4}\), independent Quality score, reference KL,
constraint residuals, accepted updates, and exact rollback counts.

The confirmatory evaluation fixes \(N=32\) for every method; no cross-budget
sensitivity result is claimed.  The group-level maximum uses all 32 candidates
and does not require a selected response.  A deployment-style
\(\arg\max_i\mu_\omega(x,Y_i)\) response is reported only as a secondary
selector diagnostic.  Held-out KL is computed on common base trajectories, so
the paired comparison measures policy motion under the same prompt protocol
rather than differences in test-set difficulty.

For completeness, the top-category diagnostic used in the cached evaluation
is
\begin{equation}
\label{eq:any-four}
U_N^{\mathrm{nom}}(\theta)
:=\mathbb E_{x,Y_{1:N}}\!\left[
1-\prod_{j=1}^{N}(1-p_{\omega,4}(x,Y_j))\right].
\end{equation}
It is the frozen-model probability that at least one candidate receives rating
four.  It is not a Gaussian tail approximation and is not the primary
estimand; the complete ordinal maximum remains the reported objective.

\subsection{Statistical analysis and success criteria}
\label{subsec:statistical-analysis}

All seeds are reported separately. The frozen analysis uses 10,000 draws of
a seed-then-prompt percentile bootstrap: sample three seeds with replacement,
then independently resample 512 prompt indices within each sampled seed.
Paired method contrasts use within-seed prompt differences before resampling.
The resulting 2.5th and 97.5th percentiles are the reported intervals. Because
prompts are shared across seeds but resampled independently within them, these
are intervals under the frozen nested-resampling convention, not an exact
crossed seed--prompt population confidence guarantee. Three seeds further
limit precision of inference over training randomness.

The recorded success gates require the robust-max difference interval's lower
endpoint to be \(>0\) against Scalar max-PPO, \(\ge-0.05\) against
Mean-PPO, and \(\ge-0.10\) against GRPO. They also require the held-out-KL
difference interval's upper endpoint to be \(<0\) against both Mean-PPO and
GRPO. The candidate nominal-mean and Quality-score interval lower endpoints
must reach their respective frozen floors. Each generated-output proxy must
have tie-adjusted preference mean \(\ge0.50\) and interval lower endpoint
\(\ge0.45\). All nine gates must pass jointly; the recorded decision is false.
These are the unchanged protocol criteria, not thresholds chosen after seeing
the results. We make no superiority claim from the positive Mean-PPO
robust-max point estimate.

Additional diagnostics average paired differences over seeds within each
prompt before sign-flip randomization. Such tests require sign-exchangeability
under their null. The historical output key \texttt{paired\_rank\_biserial}
actually records the sign-based effect \((n_+-n_-)/(n_++n_-)\), excluding
zero differences (defined as zero if all differences vanish); it does not
weight signed ranks. We retain that field name for record compatibility and
interpret it as a sign-dominance statistic. Individual intervals and auxiliary
tests are descriptive and are not simultaneous confidence guarantees across
all methods or endpoints.

\subsection{Anti-exploitation and human-grounded diagnostics}
\label{subsec:anti-exploitation}

Newly generated outputs are compared with Mean-PPO outputs by two frozen
evaluators: the Qwen2.5-14B-Instruct judge and
\texttt{RLHFlow/ArmoRM-Llama3-8B-v0.1}
\citep{zheng2023judging,wang2024interpretable}.  For each seed, the diagnostic
aligns up to 512 held-out prompt groups, retains the response selected by
$\arg\max_i\mu_\omega(x,Y_i)$ from each 32-response group, and evaluates the
matched pair.  Qwen judgments use randomized presentation order and a reverse
order check; ArmoRM supplies an independent pointwise score.  Each evaluator
must first pass its frozen RewardBench calibration gate (between 512 and 1,024
calibration pairs, with the frozen split and hashes recorded by the protocol).
Pairwise judgments are model-based proxies, not human labels.  The complete
group-level diagnostic is run separately on all 32 candidates and reports
robust/nominal maximum, probability of any rating-4 response, and the paired
group delta; selected-response diagnostics additionally report length,
character count, repeated-trigram rate, unique-word ratio, refusal behavior,
and upper-tail strata.

\begin{table}[H]
\caption{Independent generated-output proxy comparisons against Mean-PPO.
Preference is tie-adjusted, so $0.5$ denotes parity; intervals are three-seed
hierarchical-bootstrap 95\% intervals.}
\label{tab:proxy-diagnostics}
\centering\footnotesize
\begin{tabular}{lcc}
\toprule
Evaluator & Preference & 95\% CI\\
\midrule
Qwen2.5-14B judge & 0.494 & [0.478, 0.510]\\
ArmoRM-Llama3-8B-v0.1 & 0.502 & [0.458, 0.544]\\
\bottomrule
\end{tabular}
\end{table}

Both point estimates are near parity and both intervals contain 0.5. These
intervals establish neither a quality difference nor equivalence to Mean-PPO;
in particular, the Qwen proxy does not pass its declared gate.  A
target-RM increase accompanied by proxy degradation would be evidence
consistent with reward-model exploitation; these results do not show that
signature, but the model-based design cannot establish its absence.

RewardBench and Skywork chosen--rejected lockboxes contain previously collected
human labels.  For each policy, length-normalized chosen--rejected policy
log-probability accuracy and margin are reported as preference-retention
diagnostics.  They are not Best-of-32 endpoints and do not label newly sampled
policy responses.  In particular, a new response cannot be declared human
preferred merely because it was generated for a prompt that also appears in a
fixed preference benchmark.

No new human annotation is collected in this protocol.  Therefore, even if
all anti-exploitation gates pass, the permitted claim is
\emph{no evidence of exploitation under the measured model-based and
existing-label diagnostics}.  The experiment cannot establish absence of
reward hacking or direct human preference for generated outputs.  A blinded
group-rating packet is emitted so that fresh human evaluation can be added
later without regenerating or selecting outputs.

\subsection{Reproducibility and reporting}
\label{subsec:reproducibility}

Code, model adapters, data partitions, floor samples, calibration reports, and
protocol files are content-hashed before training.  A real-GPU smoke test must
execute every declared objective, verify finite applicable metrics, verify
constraint activation, and test exact parameter and optimizer rollback.
Formal training and evaluation run only after these checks pass.  Technical
failures are versioned and corrected with regression tests.  Unfavorable
scientific outcomes are reported without changing seeds, thresholds,
comparators, evaluators, or success criteria.

The 1.5B policy is evaluated by 14B reward and proxy models.  This capacity
gap may improve discrimination but may also create systematic evaluator shift.
The three-seed design measures optimization variability; it does not replace
replication at other model scales or fresh blinded labels.
